% 原书第 26—31 页；第 2.3 节的五个小节。
\section{因果模型背后的物理结构}
\label{sec:physical-structure}

本章最后简要讨论因果模型与物理学之间的联系。只关心数学结构和统计结构的读者，可以跳过这一部分。

\subsection{时间的作用}
\label{sec:ch02-role-of-time}

第 2.1 节有一个显而易见的缺口：没有讨论时间的作用。事实上，物理学通过排除从未来到过去的因果作用，将因果性纳入了基本定律。\footnote{更准确地说，一个事件只能影响位于其光锥内的事件，因为按照相对论，任何信号的传播速度都不能超过真空中的光速。}
不过，这并没有消除因果推断中的所有问题。Simon［1953］早已认识到，时间顺序虽然能够提供一种有用的不对称性，但真正重要的是不对称性，而非时间上的先后次序。

在微观层面，人们普遍认为，无论经典系统还是量子力学系统，其时间演化都是可逆的。这似乎与我们的直觉矛盾：我们认为世界的演化具有方向性，并相信如果时间倒流，我们能够察觉。对此，可以从两个角度消解矛盾。第一，假设我们拥有一种状态复杂度度量［Bennett, 1982；Zurek, 1989］，并从复杂度很低的状态出发。在这种情况下，只要时间演化具有足够的遍历性，复杂度就倾向于增加。第二，假设我们考察的是开放系统。即使封闭系统的时间演化是可逆的，例如量子力学中的幺正时间演化，与环境发生相互作用的开放子系统，其时间演化在一般情形下也未必可逆。

\subsection{物理定律}
\label{sec:ch02-physical-laws}

下面的例子有助于讨论一个经常引发争论的因果问题。理想气体定律规定，压强 $p$、体积 $V$、物质的量 $n$ 和绝对温度 $T$ 满足
\begin{equation}
 p\cdot V=n\cdot R\cdot T,
 \tag{2.2}
\end{equation}
其中，$R$ 是理想气体常数。例如，如果改变一定量气体所占据的体积 $V$，那么压强 $p$、温度 $T$，或者二者都会改变；具体如何变化，取决于干预的确切设置。反过来，如果改变 $T$，那么 $V$、$p$，或者二者都会改变。如果保持 $p$ 不变，那么至少近似地说，我们可以构造一个涉及 $T$ 和 $V$ 的因果环。究竟是什么导致了什么？有时人们据此主张，这类定律说明，除非系统显式涉及时间，否则谈论因果性就没有意义。下面我们将说明，这种看法会产生误导。气体定律 (2.2) 描述的是底层动力系统的\textbf{平衡状态（equilibrium state）}。把它写成一个简单方程，并不能充分说明哪些干预原则上可行，以及这些干预会产生什么作用。结构因果模型及其相应的有向无环图，确实能够提供这些信息。不过，对于一般的非平衡系统，给定动力系统是否能够导出结构因果模型，以及如何导出，都是困难的问题。

\subsection{带环的赋值}
\label{sec:ch02-cyclic-assignments}

我们将结构因果模型视为对随时间展开的底层过程的抽象。对于这些底层过程，反馈环本身并不构成问题：只要时间尺度足够精细，这些环就会沿时间轴展开。这里假设不存在瞬时相互作用，而光速有限，可以说已排除了这种相互作用。

虽然随时间展开的过程本身不含环，从这种过程导出的结构因果模型却可能含环。例如，使用后文注 6.5 和注 6.7 提到的方法，可以得到只涉及不再依赖时间的量的结构因果模型，而这样的模型可能具有环。在这种系统中定义一般干预会稍微困难一些，但某些类型的干预仍应可行。例如，可能可以实施将一个变量固定为某一取值的\textbf{硬干预（hard intervention）}；在底层微分方程组中加入强迫项，就可能在物理上实现这种干预，见注 6.7。这会切断环，随后便能够导出相应的干预分布。

然而，对于一组带环的结构赋值，可能无法导出相应的观测分布。考虑以下两个赋值：
\[
 \begin{aligned}
 X&\coloneqq f_X(Y,N_X),\\
 Y&\coloneqq f_Y(X,N_Y),
 \end{aligned}
\]
以及噪声变量 $N_X\indep N_Y$。与无环模型一样，我们把噪声和函数视为给定，并尝试计算它们所导出的 $X$ 与 $Y$ 的联合分布。为此，从第一个赋值 $X\coloneqq f_X(Y,N_X)$ 出发，代入某个初始的 $Y$。这样便得到一个 $X$，再将其代入另一个赋值。假设反复迭代这两个赋值后，收敛到某个不动点。该不动点就会对应于 $X,Y$ 的一个联合分布，使两个结构赋值作为随机变量之间的等式同时成立。\footnote{赋值作为随机变量之间的等式成立，意味着我们考察的是一组系统，它们的区别在于噪声变量的实现值不同。每一组噪声实现值都会导出一组可能不同的 $X,Y$ 实现值；因此，噪声的分布就导出了 $X,Y$ 上的分布。}
注意，这里假设每一步都使用同一对 $N_X,N_Y$，而不是它们相互独立的副本。

然而，$X,Y$ 的这种平衡状态并非总是存在；即使存在，也不一定能通过上述迭代找到。在线性情形下，Lacerda et al.［2008］和 Hyttinen et al.［2012］已对此作出分析，亦见 Lauritzen and Richardson［2002］。进一步细节参见注 6.5。

两个带环的结构赋值未必能够导出一个同时满足二者的分布。这一现象与我们对结构因果模型的理解相符：结构因果模型是底层物理过程的抽象，而这种抽象作为因果模型，只在有限的范围内有效。要理解一般的带环系统，可能难免要研究微分方程组，而不能只研究结构因果模型。另一方面，在某些受限情形下，停留在结构因果模型这一更表层的现象描述层次上，仍然可能有意义，例如参见 Mooij et al.［2013］。可以推测，结构因果模型或结构方程模型固有的这种困难，或许部分解释了如下历史变化：计量经济学界起初将结构方程模型视为因果模型，后来其中一些研究者却放弃了这种解释，转而仅把结构方程视为代数方程。

\subsection{干预的可行性}
\label{sec:ch02-feasibility-interventions}

我们使用机制独立原理，来说明为何能够每次只影响一个机制或一个结构赋值。真实系统虽然可能允许这种干预，但也会存在同时替换多个赋值的干预。从直观的物理意义上说，前一种干预可以看作更加基本。如果把多个基本干预组合起来，原则上可能通过精细调节，使它们相互配合；我们会认为，这违反了独立性原理 2.1 的某种形式，见原书第 24 页脚注 8。人们或许希望，“自然”的组合干预不会破坏独立性。但是，要判断一项干预在这种意义上是否“自然”，需要知道因果结构；而我们最初试图利用这类原理进行因果学习时，恰恰还不知道这一结构。归根结底，我们可以尝试诉诸物理学，判断哪些操作是基本的或自然的。

对物理系统而言，哪些操作是基本操作？这一问题在现代量子信息论中具有关键地位，并与物理相互作用结构的分析密切相关。\footnote{供有兴趣的读者参考：由 $n$ 个二能级量子系统构成的系统，可用 $2^n$ 维 Hilbert 空间 $\mathbb{C}^2\otimes\cdots\otimes\mathbb{C}^2$ 描述。作用于这一 Hilbert 空间的幺正算子对应于物理过程。对于若干这样的系统，研究者已说明如何实现一些“基本”幺正算子：它们至多作用于 $n$ 个张量分量中的两个［Nielsen and Chuang, 2000］，在其余 $n-2$ 个分量上的作用为恒等变换。此后，通过复合这些算子，就能够近似生成其他任意幺正算子［DiVincenzo, 1995］。尽管“基本”幺正操作集合显然不止这一种选择，但考虑到物理相互作用的结构，这一选择似乎是自然的。}
类似地，我们认为，理解因果关系背后的物理机制，有时可以解释为什么某些干预是自然的，而其他干预十分复杂；这种区别实质上界定了不同结构方程所代表的“模块”。

\subsection{原因与机制的独立性和热力学时间箭头}
\label{sec:ch02-thermodynamic-arrow}

下面通过讨论和一个简化模型，说明如何将机制独立原理理解为一项物理学原理。为此，考察两个变量的特殊情形，并将下述假定作为原理 2.1 的一个特例。

\begin{figure}[htbp]
\centering
\begin{tikzpicture}[causal,x=1cm,y=1cm,line cap=round]
% 三条平行入射轨迹及散射后的三条不同方向轨迹，与原图对应。
\fill (0,0) circle (.13);
\fill (.1,-.55) circle (.13);
\fill (.75,-.62) circle (.13);
\draw[thick] (0,0)--(2.15,1.92);
\draw[thick] (.1,-.55)--(2.25,1.37);
\draw[thick] (.75,-.62)--(2.9,1.30);
\draw[ed,thick] (2.15,1.92)--(.94,1.66);
\draw[ed,thick] (2.25,1.37)--(1.5,.78);
\draw[ed,thick] (2.9,1.30)--(2.31,.13);
\path[draw=gray!70,fill=gray!35]
 (2.78,2.72)--(2.98,2.02)--(3.72,2.02)--(3.13,1.59)
 --(3.36,.88)--(2.78,1.33)--(2.19,.88)--(2.42,1.59)
 --(1.83,2.02)--(2.57,2.02)--cycle;
\end{tikzpicture}
\caption{初始状态与动力学定律相互独立的简单例子：一束粒子被物体散射。出射粒子包含有关物体的信息，而入射粒子不包含这种信息。}
\label{fig:ch02-scattering}
\end{figure}


\FloatBarrier
\Needspace{5\baselineskip}
\begin{principle}[初始状态与动力学定律]
设 $s$ 是物理系统的初始状态，$M$ 是一个映射，用来描述系统动力学在某段固定时间内的演化作用，那么 $s$ 与 $M$ 相互独立。这里假定，初始状态在定义上是此前尚未与该动力学发生相互作用的状态。
\end{principle}

这里，“初始”状态 $s$ 与“最终”状态 $M(s)$ 被分别视为“原因”和“结果”。因此，$M$ 就是连接原因与结果的机制。为避免得出错误结论，需要解释原理 2.2 的最后一句。下面借助一个直观例子来说明其含义。

图 2.4 展示了一种情景：初始状态与动力学的独立性如此自然，以至于我们将其视为理所当然。一束包含 $n$ 个粒子的粒子束，最初沿着完全相同的方向朝某个物体运动，随后被该物体散射到各个方向。出射粒子的方向包含有关物体的信息，而入射粒子束不包含这种信息。“粒子最初沿完全相同方向传播”的假设当然可以放宽。即使入射束存在一定程度的无序，出射束仍然可能包含有关物体的信息。事实上，视觉和摄影之所以可能，正是因为光子携带了有关使其发生散射的物体的信息。

只要“人为设计”一束入射粒子，使所有粒子在散射后沿同一方向传播，我们就能轻易地构造出这一情景的时间反演。下面说明在这种情况下应如何理解原理 2.2。显然，只有知晓物体形状的机器或主体，才能相应地引导粒子，制备出这样的粒子束。事实上，从未与物体接触过的粒子，不可能预先包含有关该物体的信息。因此，只要把引导粒子的过程视为机制的一部分，并且不把人为设计的粒子束状态称为“初始状态”，原理 2.2 仍可成立。此时，初始状态应指粒子尚未被赋予精细调节后的动量之前的状态。

摄影图像展示的是过去已经发生的事情，而非未来将发生的事情，这是过去与未来之间最显而易见的不对称性之一。前面的讨论表明，可以将这种不对称性视为原理 2.2 的一个推论。因此，该原理将原因与结果之间的不对称性，与我们视为理所当然的过去和未来之间的不对称性联系起来。

\Needspace{8\baselineskip}
上文以非形式化方式说明了原理 2.1 与物理学中过去和未来不对称性的关系。最后简要提及，Janzing et al.［2016］已借助算法信息论，更严格地建立了这一联系。正如原理 4.13 用算法独立性来形式化 $P_C$ 与 $P_{E\mid C}$ 的独立性，原理 2.2 也可以解释为 $s$ 与 $M$ 的算法独立性。Janzing et al.［2016，定理 1］证明：对于任意双射 $M$，如果接受 Bennett［1982］与 Zurek［1989］的主张，认为 Kolmogorov 复杂度（见第 4.1.9 节）是物理熵的恰当形式化，那么原理 2.2 就意味着，$M(s)$ 的物理熵不可能小于 $s$ 的熵，但允许相差一个加法常数。因此，原理 2.2 导出了物理学中通常所说的时间箭头意义下的熵不减。
\FloatBarrier
